% =============================================================================
% Discussion
% =============================================================================
\section{Discussion}\label{sec:discussion}

\subsection{Connection to Kac's question}

Kac's original problem asked whether one could infer geometry from a spectrum.
Our answer is deliberately narrower and more practical.  For the present
viscoelastic shell model, the results suggest that one \emph{can} hear the shape
of an oblate organ, in the sense that the local inverse problem for
$(a,c,E)$ is well conditioned and statistically finite.  At the spherical
limit, one cannot: the Jacobian becomes rank deficient and the Fisher matrix
loses invertibility.

The analogy with Kac should therefore be handled with care.  We do not prove a
global uniqueness theorem of the kind established by Sylvester and
Uhlmann~\cite{SylvesterUhlmann1987} for the Calder\'on inverse conductivity
problem, nor do we exclude the possibility of remote isospectral parameter sets.
Our claim is about \emph{practical identifiability} in a finite-dimensional
model relevant to measurement and inversion.  That is a less grand claim, but
also the one an experimentalist can actually use.

\subsection{Mechanism: mode-dependent curvature sampling}

The analytical results in Section~\ref{sec:theory} clarify why oblate geometry
restores identifiability.  Proposition~\ref{prop:chain_rule} shows that \emph{any}
model built on $R_\mathrm{eq} = (a^2 c)^{1/3}$ alone has a rank-deficient
Jacobian, because the chain rule forces the $a$- and $c$-columns to be exactly
proportional for all modes.  Proposition~\ref{prop:eccentricity} then shows
that the Ritz model, which integrates strain energy over the actual spheroidal
surface, introduces mode-dependent corrections at order $\varepsilon^2$.

The mechanism is \emph{mode-dependent curvature sampling}.  Different Legendre
modes $P_n(\eta)$ weight the meridional curvature gradient differently: low-order
modes sample both poles and equator broadly, while higher-order modes
concentrate energy near the equator where the curvature gradient is steepest.
This differential sampling breaks the proportionality of the $a$- and $c$-columns
and restores full rank.  The condition number diverges as
$\kappa \sim C\,\varepsilon^{-\alpha}$ as $\varepsilon \to 0$, with the
exponent $\alpha$ confirmed numerically in Section~\ref{sec:results}.

This interpretation is straightforward and does not invoke group-theoretic
symmetry breaking.  The identifiability transition is not a consequence of
reducing the symmetry group from $\mathrm{SO}(3)$ to $\mathrm{SO}(2)$; it is a
consequence of mode-dependent curvature weighting in the variational
integrals.  The distinction matters because the curvature-sampling explanation
is constructive --- it identifies the specific integrals responsible --- whereas
a symmetry-based framing would be descriptive at best and misleading at worst.

\subsection{Implications for vibroacoustic elastography}

The results are encouraging for vibroacoustic elastography and related spectral
methods.  A small number of low-order flexural resonances may contain enough
information to estimate both geometry and stiffness, provided that the operating
point is sufficiently aspherical and provided that enough modes are measured.
This is attractive because resonant frequencies are often easier to estimate
robustly than full-field mode shapes or transient wavefields.

The Cram\'er--Rao bounds give the result practical content.  At $1\%$ frequency
noise, the uncertainty in $c$ is modest, while $a$ and $E$ remain more weakly
resolved but still finite.  That anisotropy matters for experimental design:
one should not expect the same confidence in all parameters, even when the
Jacobian has full rank.

\subsection{Why model choice matters}

The most important modelling lesson is that forward-model convenience and
inverse-model usefulness are not the same thing.  A spherical reduction may be
acceptable for a rough forward estimate of resonance frequency, but it is a poor
foundation for parameter recovery.  Spherical simplifications are serviceable on
the blackboard; in inverse problems they become noticeably less cooperative.

At the same time, the present study does \emph{not} rely on a full oblate shell
theory.  The model analysed here is still an equivalent-sphere approximation.
That limitation should be kept in view when interpreting the phrase ``hear the
shape''.  The present results show what is identifiable within the implemented
model, not what Nature has promised in advance.

\subsection{Relation to existing inverse methods}

Inverse vibration problems have long been studied in structural
dynamics~\cite{Gladwell2004,Gladwell2005}, and model-based elastography has
developed a wide range of optimisation and regularisation
strategies~\cite{Doyley2012}.  In acoustics and biomechanics, several studies
have estimated material parameters by fitting dynamic response models to
spectral data~\cite{Aernouts2012}.  The present contribution is more modest but
useful: before attempting a sophisticated inversion, we ask whether the chosen
measurements can, in principle, resolve the chosen parameters with acceptable
conditioning.

That perspective is valuable because poor inverse performance is often blamed on
optimisers, priors, or noise models when the underlying issue is simpler: the
Jacobian is nearly singular.  A condition-number and Fisher-information analysis
cannot solve the inverse problem by itself, but it can prevent one from asking
the data for information they do not contain.

\subsection{Local versus global identifiability}

The round-trip tests demonstrate very accurate local recovery, but global
identifiability is a separate matter.  The three-mode system is especially
vulnerable because it is exactly determined.  Local minima can occur, and a
small residual does not necessarily imply the correct parameter set when the
search is started far away.  This is why the five-mode result matters more for
practice than the machine-precision local recovery.

Whether the spectral map is globally injective --- that is, whether the oblate
shell inverse problem admits a unique solution over the full parameter domain
--- remains an open question connecting to the spectral geometry programme
initiated by Kac~\cite{Kac1966} and the perturbation theory of
Kato~\cite{Kato1966}.  The eigenvalue perturbation results of
Hald~\cite{Hald1978,Hald1984CPAM} and the shape optimisation framework of
Henrot~\cite{Henrot2006} provide relevant tools, but a proof of global
uniqueness for the present coupled fluid--shell system would require substantial
additional analysis.

Our recommendation is therefore not based on elegance but on behaviour:
\begin{itemize}
  \item use three modes only for local studies or controlled round-trip tests;
  \item use at least five flexural modes for robust inversion in realistic
  settings.
\end{itemize}

\subsection{Limitations}

Several limitations remain.
\begin{enumerate}
  \item The forward model is an equivalent-sphere approximation.  It does not
  contain explicit $\zeta_g$-dependent shell terms, so the present paper should
  not be read as a definitive study of full spheroidal shell identifiability.
  \item The analysis is local.  We do not prove global injectivity of the map
  $(a,c,E)\mapsto \{f_n\}$.
  \item Synthetic inversion uses the same model for data generation and
  inversion.  This isolates identifiability, but it avoids model-form error.
  \item The mode numbers are assumed known.  In experiment, mode assignment may
  itself be uncertain.
  \item Only three parameters are estimated.  Allowing $h$, $\nu$, or fluid
  properties to vary would enlarge the parameter space and may degrade
  conditioning.
\end{enumerate}

\subsection{Future directions}

The natural next step is a genuine Ritz or finite-element oblate shell model
with explicit aspect-ratio dependence, followed by the same identifiability
analysis.  Shape-calculus tools~\cite{DelfourZolesio2011} would allow
Hadamard-type Fr\'echet derivatives of the eigenvalues with respect to the
domain boundary, enabling a geometry-aware sensitivity analysis that goes beyond
the parametric approach taken here.  That would show whether the encouraging
results reported here strengthen further when the geometry is represented more
faithfully.  The second
priority is experimental validation, ideally first on phantoms of known
geometry and stiffness and then on biological specimens.  A Bayesian treatment,
with priors on tissue properties and uncertainty in mode identification, would
also be a useful extension.

The results suggest that geometric asymmetry plays a fundamental role in
practical identifiability of coupled fluid--shell systems, a connection that
merits further investigation using the tools of shape calculus and spectral
perturbation theory~\cite{Kato1966,Henrot2006}.

\subsection{Open Questions and Conjectures}\label{sec:open_questions}

Propositions~\ref{prop:chain_rule} and~\ref{prop:eccentricity} identify the
mechanism by which the spherical degeneracy is lifted, but they stop short of a
complete quantitative theory.  To make the outstanding mathematical programme
explicit, we record three conjectures suggested by the numerical evidence.
These are stated for the genuine oblate Ritz model, since that is the setting
in which a proof would matter.  Let
$\mathcal{N}=\{n_1,\ldots,n_N\}$ denote the measured flexural mode set, let
$\varepsilon=\sqrt{1-c^2/a^2}$ be the eccentricity, and let
$\mathbf{J}_s(\varepsilon;\mathcal{N})$ denote the corresponding scaled
Jacobian of the frequency--parameter map.

\begin{conjecture}[Rate of identifiability restoration]
\label{conj:power_law}
For the Ritz model with fixed mode set $\mathcal{N}$, there exist constants
$C(\mathcal{N})>0$, $\alpha(\mathcal{N})>0$, and $\varepsilon_0>0$ such that
\begin{equation}
  \kappa\bigl(\mathbf{J}_s(\varepsilon;\mathcal{N})\bigr)
  =
  C(\mathcal{N})\,\varepsilon^{-\alpha(\mathcal{N})}\bigl(1+o(1)\bigr)
  \qquad \text{as } \varepsilon \to 0^+,
\end{equation}
where the exponent $\alpha(\mathcal{N})$ depends only on the number of modes
$N$ and on the selected mode set $\mathcal{N}$.
\end{conjecture}

\begin{conjecture}[Minimum modes for identifiability]
\label{conj:min_modes}
Consider the four-parameter material map with parameter vector
\begin{equation}
  \boldsymbol{\vartheta} = (E,h,\rho_w,\rho_f)^\top,
\end{equation}
at fixed geometry and fixed $(\nu,K_f,P_\mathrm{iap})$, reduced by the
one-dimensional similarity scaling identified by the Buckingham--$\Pi$ analysis
so that the effective parameter dimension is three.  Then three measured
flexural modes suffice for local identifiability if and only if
$\varepsilon>0$: equivalently, the reduced scaled Jacobian has full column rank
for some three-mode set when $\varepsilon>0$, whereas no finite mode set yields
full column rank when $\varepsilon=0$.
\end{conjecture}

\begin{conjecture}[Universality of scaling]
\label{conj:universality}
For the nested flexural mode family $\mathcal{N}_N=\{2,\ldots,N+1\}$, the
Buckingham--$\Pi$ groups induced by the Ritz equations are complete in the sense
that there exists a universal function $\mathcal{K}$ satisfying
\begin{equation}
  \kappa\bigl(\mathbf{J}_s(\varepsilon;\mathcal{N}_N)\bigr)
  =
  \mathcal{K}(\varepsilon,N),
\end{equation}
independently of the absolute values of
$a$, $c$, $E$, $h$, $\rho_w$, and $\rho_f$, provided the dimensionless groups
other than $\varepsilon$ are held fixed.
\end{conjecture}

Conjecture~\ref{conj:power_law} is the most tantalising of the three.  The
numerics already suggest a clean power law, but a proof would turn a useful
plot into an asymptotic theorem.  That would be the crown jewel of the present
identifiability programme; at the time of writing, it remains work in progress.
